\section{A direct coarse bound from algebraic \texorpdfstring{$K$}{K}-theory}
\label{kt-section}

The coarse cohomology estimate has a shorter route through established
comparison theorems in algebraic \(K\)-theory.  We state these inputs
precisely, including the bound at the prime \(2\), and prove all the
resulting order estimates.  This gives an alternative to
assertion~\ref{sc-few} for the coarse estimate; that assertion itself
is proved independently in Theorem~\ref{sw-few}.

Let \(L\) be a totally imaginary number field, let
\[
 n=[L:\mathbb Q],\qquad D=|\operatorname{disc}L|,\qquad
 h=\#\operatorname{Cl}(L),
\]
and let \(p\) be a prime such that \(\mu_p\subset L\).
For \(p=2\) this last condition is automatic.
Let \(S\) be the finite primes above \(p\), let
\(R=\mathcal O_L[1/p]\), and let \(G_{L,S}\) be the Galois
group of the maximal algebraic extension unramified outside \(S\).
All \(p\)-adic cohomology groups below use continuous cohomology.
For a finite abelian group \(A\), its \(p\)-primary subgroup is
\(A\{p\}=\{a:p^r a=0\text{ for some }r\ge0\}\).

\subsection{The precise classical inputs}

Write \(\operatorname{GL}(\mathcal O_L)\) for the union of the
groups \(\operatorname{GL}_N(\mathcal O_L)\) under the usual
upper-left inclusions.  Quillen's construction gives a connected
infinite loop space
\(X=B\operatorname{GL}(\mathcal O_L)^+\).
It is the base component of the \(K\)-theory infinite
loop space, equivalently the infinite loop space of the
connected cover of its spectrum; taking this cover
preserves all positive homotopy groups.
Here the plus construction preserves integral homology, and
\[
 K_j(\mathcal O_L)=\pi_j(X)\quad(j\ge1),\qquad
 H_j(X,\mathbb Z)=
 H_j(\operatorname{GL}(\mathcal O_L),\mathbb Z).
 \tag{KT1}
\]
An infinite loop space is a pointed space equipped with
compatible deloopings in all positive degrees.  We use this
structure only to apply the Hurewicz theorem with an exponent
bound stated below.

The required results are the following.

\begin{enumerate}
\item Quillen's finite-generation theorem and Borel's rank
theorem imply that \(K_{2E}(\mathcal O_L)\) is finite for
every \(E\ge1\).
\cite[IV, Theorems 6.9 and 1.17--1.18]{weibel-kbook}
\item Homological stability gives
\[
 H_m(\operatorname{GL}_N(\mathcal O_L),\mathbb Z)
       \simeq H_m(\operatorname{GL}(\mathcal O_L),\mathbb Z)
       \quad\text{if }N\ge\max(2m+1,m+2).
 \tag{KT2}
\]
This is the stable-range form of the usual integral
stability theorem.\cite[IV, Corollary 1.14.1]{weibel-kbook}
\item The Quillen--Lichtenbaum comparison, valid also at \(2\)
for totally imaginary \(L\), identifies
\[
 K_{2E}(\mathcal O_L)\{p\}
       \simeq H^2_{\mathrm{\acute et}}
                  (\operatorname{Spec}R,\mathbb Z_p(E+1)),
                    \qquad E\ge1.
 \tag{KT3}
\]
\cite[VI, Theorem 8.2]{weibel-kbook}
For \(p\)-power torsion coefficients, the usual
\'etale--Galois comparison identifies the right-hand
cohomology with that of \(G_{L,S}\), since \(p\) is
invertible on \(\operatorname{Spec}R\).
It also identifies their compatible inverse limits for
\(\mathbb Z_p(E+1)\).
\cite[II, Proposition 2.9]{milne-adt}
\item If \(Y\) is a connected infinite loop space and \(j\ge2\),
the \(p\)-primary part of the kernel of
\[
            \pi_j(Y)\longrightarrow H_j(Y,\mathbb Z)
\]
is killed by \(p^{b_p(j)}\), where one may take
\[
 b_2(j)=\left\lceil\frac{j-1}{2}\right\rceil+3,\qquad
 b_p(j)=\left\lceil\frac{j+2}{2p-2}\right\rceil+1
                  \quad(p>2).
 \tag{KT4}
\]
This is Mathew's exponent theorem with connectivity
parameter \(m=1\).\cite[Theorem 6.3]{mathew-hurewicz}
In particular \(b_p(2E)\le E+3\) for every prime \(p\).
\end{enumerate}

These theorems are used as established inputs; their general
proofs are not reproduced here.  In (KT3), the integral
\(p\)-primary group is the one over \(\mathcal O_L\), not
over the field \(L\).  The comparison includes the
localization at \(p\) displayed on the right.
For the \(p\)-adic comparison, the finite-coefficient
\(H^1\) groups are finite.  One can see this directly:
after passing to \(L(\mu_{p^a})\), the module
\(\mathbb Z/p^a(E+1)\) is isomorphic to \(\mu_{p^a}\).
Kummer theory expresses its \(H^1\) as an extension of
the \(S\)-units modulo \(p^a\)-th powers by a subgroup
of the finite ideal class group.
The \(S\)-unit group is finitely generated: its valuation
map has kernel the ordinary units and image in
\(\mathbb Z^S\), and the ordinary unit theorem was
proved in Theorem~\ref{su-unit-theorem}.
Restriction has finite kernel, since that kernel is
cohomology of a finite group with finite coefficients.
This proves the required finiteness before base change as well.
Inverse systems of finite groups satisfy the
Mittag--Leffler condition, so the possible
\(\varprojlim^1H^1\) term vanishes and the two ordinary
inverse limits agree.

For completeness, the stable-range bound used in (KT2)
does apply to every ring of integers.  The needed ring
condition says that a unimodular triple \((a,b,c)\) can
be shortened to a unimodular pair \((a+cx,b+cy)\).
If \(c=0\), take \(x=y=0\).
Otherwise choose \(x\in\mathcal O_L\) so that \(d=a+cx\ne0\).
There are only finitely many prime ideals dividing \((d)\).
At any such prime dividing \(c\), unimodularity forces
\(b\) to be nonzero modulo that prime.
At each of the other primes, choose \(y\) modulo the prime
so that \(b+cy\ne0\).  The Chinese remainder theorem makes
all these choices simultaneous.
No prime ideal divides both \(d\) and \(b+cy\), proving
the ring condition, or stable range at most \(2\).
Thus the general bound
\(N\ge\max(2m+1,m+\operatorname{sr}(\mathcal O_L))\)
gives (KT2).

\subsection{The number of cyclic \texorpdfstring{$p$}{p}-primary factors}

We also use the classical Brauer--Hasse--Noether theorem
in the following form for the ring \(R\):
\[
 0\longrightarrow\operatorname{Br}(R)
   \longrightarrow\bigoplus_{v\in S}\mathbb Q/\mathbb Z
   \xrightarrow{\ \sum_v\ }\mathbb Q/\mathbb Z
   \longrightarrow0.
 \tag{KT5}
\]
The field is totally imaginary, so there are no real-place
terms.  Here \(\operatorname{Br}(R)=
H^2_{\mathrm{\acute et}}(\operatorname{Spec}R,\mathbb G_m)\);
the maps are the local invariants.
\cite[VI, equation (8.1.1)]{weibel-kbook}
Only the resulting elementary \(p\)-rank bound will be needed.

\begin{lemma}\label{kt-rank}
Put \(M_E=H^2(G_{L,S},\mathbb Z_p(E+1))\).
This is a finite \(p\)-group, and
\[
 \dim_{\mathbb F_p}(M_E/pM_E)
       \le |S|-1+\log_p h
       \le n+\log_p h .
 \tag{KT6}
\]
\end{lemma}
\begin{proof}
Finiteness follows from the first and third classical
inputs.  The coefficient sequence for multiplication by
\(p\) gives an injection
\[
 M_E/pM_E\ \lhook\joinrel\longrightarrow\
       H^2(G_{L,S},\mathbb F_p(E+1)).
\]
Since \(\mu_p\subset L\), all these finite twists have
trivial Galois action.  Choosing a primitive root
identifies the finite coefficient module with \(\mu_p\).
Use the \'etale--Galois comparison to compute the
last group over \(\operatorname{Spec}R\).

The Kummer sequence
\[
 1\longrightarrow\mu_p\longrightarrow\mathbb G_m
       \xrightarrow{\,(\cdot)^p\,}\mathbb G_m
       \longrightarrow1
\]
is exact on the \'etale site: a unit acquires a \(p\)-th
root in a finite \'etale algebra because \(p\) is invertible.
Its cohomology sequence, with
\(H^1(\mathbb G_m)=\operatorname{Pic}(R)\), gives
\[
 0\longrightarrow\operatorname{Pic}(R)/p
  \longrightarrow H^2_{\mathrm{\acute et}}(R,\mu_p)
  \longrightarrow\operatorname{Br}(R)[p]
  \longrightarrow0.
 \tag{KT7}
\]
Every fractional ideal of \(R\) comes from an ideal of
\(\mathcal O_L\); factorization into prime ideals shows
that localization merely kills the classes of the
primes in \(S\).  Thus \(\operatorname{Pic}(R)\) is a
quotient of \(\operatorname{Cl}(L)\), and its quotient
by \(p\) has at most \(h\) elements.

Equation (KT5) identifies \(\operatorname{Br}(R)[p]\)
with the kernel of the sum map
\(\mathbb F_p^{S}\to\mathbb F_p\).
There is at least one prime over \(p\), so its dimension
is exactly \(|S|-1\).  Add the dimensions in (KT7).
Finally \(|S|\le n\), since
\(\sum_{v\mid p}e_vf_v=n\) and each summand is positive.
This proves (KT6).
\end{proof}

\subsection{The complete order comparison}

\begin{proposition}\label{kt-additive}
For every \(E\ge1\),
\[
 \log\#H^2(G_{L,S},\mathbb Z_p(E+1))
 \le
 \log t_L(4E+6,2E)
       +(E+3)\bigl(n\log p+\log h\bigr).
 \tag{KT8}
\]
\end{proposition}
\begin{proof}
Let \(A=K_{2E}(\mathcal O_L)\{p\}\), which is finite.
The ordinary Hurewicz map and (KT1) give
\[
 f:A\longrightarrow
               H_{2E}(\operatorname{GL}(\mathcal O_L),\mathbb Z).
\]
The image of a finite group lies in the torsion subgroup.
For \(m=2E\), the choice \(N=4E+6\) satisfies (KT2),
so
\[
                  |\operatorname{im}f|
                              \le t_L(4E+6,2E).
\]
This is where finiteness of the even \(K\)-group is
needed; it is supplied explicitly above.

By (KT4), \(\ker f\) is killed by \(p^{E+3}\).
Write \(r=\dim_{\mathbb F_p}(A/pA)\).
For any finite abelian \(p\)-group the kernel and
cokernel of multiplication by \(p\) have equal orders,
so \(|A[p]|=p^r\).
For \(j\ge1\), multiplication by \(p\) maps
\(A[p^{j+1}]\) to \(A[p^j]\) and has kernel \(A[p]\).
Induction therefore gives
\[
                 |A[p^j]|\le p^{jr}.
\]
Consequently \(|\ker f|\le p^{(E+3)r}\).
By (KT3) and Lemma~\ref{kt-rank},
\(r\le n+\log_p h\).
The identity \(|A|=|\ker f|\,|\operatorname{im}f|\)
and (KT3) now prove (KT8).
\end{proof}

\begin{theorem}\label{kt-fourteen}
For every totally imaginary \(L\), every prime \(p\)
with \(\mu_p\subset L\), and every \(E\ge1\), one has
\[
 \log\#H^2(G_{L,S},\mathbb Z_p(E+1))
       \le\exp\bigl((2En)^{14}\log(2D)\bigr).
 \tag{KT9}
\]
With the stated global duality input, the strict Selmer
group with coefficient module
\(\mathbb Q_p/\mathbb Z_p(-E)\) is finite and also satisfies
\[
        \log\#\mathcal S_p(L,E)
            \le\exp\bigl((2En)^{14}\log(2D)\bigr).
 \tag{KT10}
\]
\end{theorem}
\begin{proof}
Put \(x=2En\ge4\) and \(Y=\log(2D)\ge\log2>1/2\).
The arithmetic complex and
Proposition~\ref{ht-arithmetic-numerical} give
\[
                 \log t_L(4E+6,2E)\le e^{x^{13}Y}.
\]
Corollary~\ref{sl-class-number}, together with
\(\log(1+t)\le t\) for \(t\ge0\), gives
\[
 \log h\le\tfrac12\log D
        +(n-1)\log(1+\tfrac12\log D)
            \le\tfrac n2\log D\le\tfrac n2Y.
\]
The inclusion \(\mu_p\subset L\) implies \(p-1\le n\).
For odd \(p\), this is the degree of the prime cyclotomic
field, proved by the Eisenstein calculation in
Section~\ref{sl-section}; it also holds for \(p=2\).
Thus \(\log p\le\log(n+1)\le n\).
Since \(E+3\le4E\) and \(n\le x/2\), the extra term in
(KT8) is at most
\[
 4E(n^2+nY/2)=2xn+xY
       \le x^2+xY\le3x^2Y\le x^{13}Y.
\]
Here \(Y>1/2\) and \(x\ge4\) give the penultimate
inequality, and \(x^{11}\ge3\) gives the last.
Therefore (KT8) is at most
\[
 e^{x^{13}Y}+x^{13}Y
       \le2e^{x^{13}Y}\le e^{x^{14}Y}.
\]
For the last step, use
\((x-1)x^{13}Y\ge\log2\).
This proves (KT9), with all small primes included.

The global duality statement before
Proposition~\ref{ht-h2-from-selmer} identifies
\(\mathcal S_p(L,E)^\vee\) with a subgroup of this
finite \(H^2\).  It follows that the Selmer group is
finite and has order no larger than that of \(H^2\),
proving (KT10).
\end{proof}

Theorem~\ref{kt-fourteen} proves the coarse cohomology
input using the explicitly cited classical theorems.
It does not use the source-specific finite-field
homology construction in assertion~\ref{sc-few}.
The independent proof of that construction is given in
Theorem~\ref{sw-few}.  The other arithmetic and analytic
constructions needed for the claimed Riemann hypothesis
remain distinct tasks.

% Shared ordinary-mathematics reductions for the RH and Dirichlet GRH routes.
\subsection{Finiteness before quantitative estimates}

Let \(F\) be a totally imaginary number field, \(p\) a prime,
and \(e\ge1\).  Let \(S\) consist of the places above \(p\)
and the archimedean places, and let \(G_{F,S}\) be the Galois
group of the maximal extension unramified outside \(S\).  Define
\[
 \mathcal S(F,p,e)=\ker\left(
 H^1(G_{F,S},\mathbf Q_p/\mathbf Z_p(-e))
 \longrightarrow\prod_{v\mid p}
             H^1(F_v,\mathbf Q_p/\mathbf Z_p(-e))\right).
\]
Thus the local condition is zero in the full local cohomology group.

\begin{proposition}\label{shared-selmer-finite}
There is an injection
\[
 \mathcal S(F,p,e)^\vee\hookrightarrow K_{2e}(\mathcal O_F)\{p\}.
\]
In particular \(\mathcal S(F,p,e)\) is finite for every prime,
including \(p=2\), and
\[
 \operatorname{length}_{\mathbf Z_p}\mathcal S(F,p,e)
       \le v_p\bigl(\#K_{2e}(\mathcal O_F)\bigr).
\]
\end{proposition}
\begin{proof}
Quillen's finite-generation theorem and Borel's rank theorem make
\(K_{2e}(\mathcal O_F)\) finite.  The Quillen--Lichtenbaum
comparison identifies its \(p\)-primary subgroup with
\(H^2_{\mathrm{\acute et}}(\mathcal O_F[1/p],\mathbf Z_p(e+1))\).
This comparison includes \(e=1\); at \(p=2\) it requires only
that \(F\) be totally imaginary, and does not require
\(\mu_p\subset F\).
\cite[IV, Theorems 1.17--1.18 and 6.9; VI, Theorem 8.2]{weibel-kbook}

Put \(A_n=(\mathbf Q_p/\mathbf Z_p(-e))[p^n]\).
Its finite Tate dual is \(T_n=\mathbf Z/p^n(e+1)\).
Write \(\mathcal S^i_S\) for the kernel of localization in
degree \(i\).  Finite-coefficient global duality gives compatible
perfect pairings between \(\mathcal S^1_S(A_n)\) and
\(\mathcal S^2_S(T_n)\).\cite[I, Theorem 4.10]{milne-adt}
Continuous cochains with discrete coefficients commute with the
union of the \(A_n\), since their compact domains have finite image.
The set \(S\) is finite, so kernels and its product commute with
this filtered colimit.  Taking duals and then inverse limits gives
\[
 \mathcal S(F,p,e)^\vee
   \simeq\varprojlim_n\mathcal S^2_S(T_n)
   \hookrightarrow\varprojlim_n H^2(G_{F,S},T_n).
\]
The finite-level \(H^1\) groups are finite, so the possible
\(\varprojlim^1H^1\) term vanishes.  The \'etale--Galois comparison
therefore identifies the last inverse limit with the integral
cohomology group above.\cite[I, Corollary 4.15; II, Proposition 2.9]{milne-adt}
The injection follows.  A discrete abelian group with finite
Pontryagin dual is finite, and has the same order as that dual.
\end{proof}

\begin{corollary}\label{shared-selmer-finite-support}
For fixed \(F\) and integer \(t\ge2\), put
\(N_{F,t}=\prod_{e=1}^{t-1}\#K_{2e}(\mathcal O_F)\).
If \(p\nmid N_{F,t}\), then \(\mathcal S(F,p,e)=0\) for
every \(1\le e<t\).  Moreover the following sum has finite support:
\[
 \sum_p(\log p)\sum_{e=1}^{t-1}
        \operatorname{length}_{\mathbf Z_p}\mathcal S(F,p,e)
                  \le\log N_{F,t}.
\]
\end{corollary}
\begin{proof}
Sum the prime-valuation bound of Proposition~\ref{shared-selmer-finite}
over the finite product defining \(N_{F,t}\).
\end{proof}
This supplies qualitative finiteness and convergence for each datum.
A useful bound uniform in the changing field and twist range still
requires a quantitative estimate.

\subsection{Fewer characters and a joint count of primes and twists}

Suppose the target character has order \(o\), and every auxiliary
character \(\psi\) has order dividing \(2^h o\) for some integer
\(h\ge0\).  If \(o\) is a power of two, a nontrivial \(\psi\)
has even order: an odd divisor of a power of two is one.  Consequently
the part of a character estimate restricted to nontrivial odd-order
auxiliary characters is empty.  This applies to the trivial target
used for RH.  It leaves the trivial character, the even-order
characters and the exceptional coefficient primes.

\begin{proposition}\label{shared-prime-twist-count}
Let \(t\ge2\) be an integer, \(m=t-1\), \(H_m=\sum_{j=1}^m1/j\), and
\(N\ge1\).  For each odd prime \(p\) in a finite set and each
\(1\le e<t\), consider at most \(N\) characters with nonnegative
lengths \(\ell(p,e,\psi)\).  Suppose that, for each \(e\), the
sum of \((\log p)\ell(p,e,\psi)\) over the terms with
\(p-1\nmid2e\) is at most \((Nt)^{C_I}\), and every weighted term
\((\log p)\ell(p,e,\psi)\) with \(p-1\mid2e\) is at most
\((pNe)^C\), where \(C_I,C\ge0\).
Then the total weighted sum is at most
\[
 m(Nt)^{C_I}+NmH_m\bigl((2m+1)Nm\bigr)^C
       \le(1+2^C)(Nt)^{\max\{C_I+1,\,2C+2\}}.
\]
\end{proposition}
\begin{proof}
Put \(j=(p-1)/2\).  The condition is \(j\mid e\), so the
number of prime--twist pairs is at most
\(\sum_{j=1}^m\lfloor m/j\rfloor\le mH_m\).
Such pairs have \(p\le2m+1\) and \(e\le m\).
Sum the stated bounds and use the character count.
For the last inequality use \(m\le t\), \(H_m\le t\),
\(pNe\le2Nt^2\), and \(N^{C+1}\le N^{2C+2}\).
\end{proof}
The exponent \(C\) must cover every remaining resonant character,
including the trivial one unless its contribution separately vanishes.
These are conditional character estimates, not consequences of finiteness.

\subsection{Retain only the actual twisted places}

Let \(K\) be the field normalizing the finite-place weights, and
let \(\mathcal W\) be the fixed finite set of twisted places.
Write \(e_w,f_w\) for the absolute ramification and inertia degrees,
\(S_0\) for the residue characteristics, and set
\[
 a_w=\frac{e_w\log Nw}{[K:\mathbf Q]},\qquad
 \kappa_p=\frac{\sum_{w\in\mathcal W,\,w\mid p}e_wf_w}
                  {[K:\mathbf Q]}\le1.
\]
Thus \(\sum_{w\in\mathcal W,\,w\mid p}a_w=\kappa_p\log p\).
Fix a set \(G\) of generator primes.  Let \(V_p\ge0\) be the
Selmer-length aggregate over the chosen twists of the same host
field, which may differ from \(K\).

\begin{proposition}\label{shared-twisted-selmer-budget}
Let \(d_j\ge0\), \(R_p\ge0\), and \(\eta_j(w)\ge0\).
Assume, for \(p\notin G\), that
\(\eta_j(w)\le d_j(R_p+V_p)\) at every selected place above \(p\).
At \(p\in G\), assume the aggregate estimate
\[
 \sum_{w\in\mathcal W,\,w\mid p}a_w\eta_j(w)
               \le d_j(\log p)(R_p+V_p).
\]
If \(\sum_{p\in S_0}(\log p)R_p\le A\), then
\[
 \sum_{w\in\mathcal W}a_w\eta_j(w)
 \le d_j\left(A+\sum_{p\in S_0}\widehat\kappa_p(\log p)V_p\right),
 \qquad
 \widehat\kappa_p=
 \begin{cases}1,&p\in G,\\ \kappa_p,&p\notin G.\end{cases}
\]
\end{proposition}
\begin{proof}
Sum the per-place bounds with their actual weights outside \(G\),
and use the given aggregate bound on \(G\).  Retain
\(\widehat\kappa_p\) on \(V_p\); on \(R_p\) use
\(0\le\widehat\kappa_p\le1\).  Every sum is finite.
\end{proof}

In particular, let \(y=x^\varepsilon\ge1\), \(b=e^y\),
\(f,m\ge1\), \(B_0=fmby\), and \(N_c\ge0\).
If \(d_j=D_0\), \(A\le C_Nfmby^2\), and
\[
 \sum_{p\in S_0}\widehat\kappa_p(\log p)V_p
                       \le C_Vy(B_0+N_c),
\]
with nonnegative uniform constants, then the congruence loss is at
most \(D_0(C_N+C_V)y(B_0+N_c)\).  This follows from
\(fmby^2=yB_0\).  The support, integral modules and local estimates
must belong to the same chosen datum.  They cannot be selected anew
to discard large lengths.  Proposition~\ref{shared-selmer-finite}
supplies finiteness; the useful uniform weighted estimate and the
integral local comparisons remain construction requirements.

